\begin{problem}[25 pts]
\statementfigure{%
It is one-dimensional motion along the x-axis. The initial position of
the object is $x(0)=\qty{15}{\metre}$. Figure on the right shows the
change in the acceleration with time. What is the x-coordinate of the
object at time $t=\qty{6}{\second}$, if the initial velocity equals to
$\qty{-10}{\metre\per\second}$?
}{\resizebox{\linewidth}{!}{\begin{tikzpicture}
\begin{axis}[width=10.5cm,height=4.5cm,axis lines=middle,
 xmin=0,xmax=11,ymin=-0.8,ymax=9.5,
 xlabel={$t$ (\unit{s})},ylabel={$a$ (\unit{\metre\per\second\squared})},
 xtick={0,2,4,6,8,10},ytick={0,2,4,6,8},
 tick label style={font=\large},label style={font=\large},
 ylabel style={at={(axis description cs:0,1)},anchor=south west,
 rotate=0,yshift=3pt},
 grid=major,grid style={PhysicsLine},clip=false]
\addplot[physics line] coordinates
 {(0,0)(2,0)(2,4)(4,4)(4,0)(8,0)(10,8)};
\draw[dashed,PhysicsBlue] (axis cs:10,0)--(axis cs:10,8)--(axis cs:0,8);
\end{axis}
\end{tikzpicture}
}}
\end{problem}
\begin{solution}
\qstep{1}{Split the motion into three equal time intervals.}
Use 0--2, 2--4 and 4--6 s, each of duration $\Delta t$.
The accelerations are $a_1=0$, $a_2$ and $a_3=0$; the index labels the interval.
The ramp after 8 s cannot affect the position at 6 s.
\qstep{2}{Write the displacement in each interval.}
For constant acceleration,
$\Delta x=v_{\rm start}\Delta t+\frac{a(\Delta t)^2}{2}$ and
$v_{\rm end}=v_{\rm start}+a\Delta t$. Therefore
\[
\begin{aligned}
\Delta x_1&=v_0\Delta t,\qquad
\Delta x_2=v_0\Delta t+\frac{a_2(\Delta t)^2}{2},\\
\Delta x_3&=(v_0+a_2\Delta t)\Delta t.
\end{aligned}
\]
\qstep{3}{Add the displacements to the initial position.}
Before substituting the data, combine the three terms:
\[
x=x_0+\Delta x_1+\Delta x_2+\Delta x_3
=x_0+3v_0\Delta t+\frac{3a_2(\Delta t)^2}{2}.
\]
\qstep{4}{Calculate the coordinate and check the graph.}
With $x_0=15$, $v_0=-10$, $\Delta t=2$ and $a_2=4$ in SI units,
\[
x=15+3(-10)(2)+\frac{3(4)(2)^2}{2}
=15-60+24=-\qty{21}{\metre}.
\]
\sourcefig{0.60\linewidth}{velocity-v2}
The areas are $-20$, $-12$ and $-\qty{4}{\metre}$, totaling
$-\qty{36}{\metre}$. This is displacement; adding the initial
\qty{15}{\metre} gives the coordinate $-\qty{21}{\metre}$.
\end{solution}

\quiznextproblem
\begin{problem}[25 pts]
The tram starts moving with the constant acceleration $a$.
The passenger waits for \qty{10}{\second} and then catches up the tram
by running after with the constant velocity of
\qty{12}{\metre\per\second}. What is the maximum acceleration of the
tram $a_{\max}$ to make this happen?
\end{problem}
\begin{solution}
\qstep{1}{Choose the model and sketch the boundary.}
Assume a shared starting position and a tram initially at rest.
The passenger waits there for time $t_{\rm w}$ and then runs at constant speed $v_{\rm P}$.
A limiting meeting counts as catching up.
\sourcefig{0.48\linewidth}{tram-limit}
The limiting line touches the tram's parabola.
\qstep{2}{Write the meeting equation.}
For $t\geq t_{\rm w}$,
\[
x_T=\frac{at^2}{2},\qquad x_P=v_{\rm P}(t-t_{\rm w}),\qquad
at^2-2v_{\rm P}t+2v_{\rm P}t_{\rm w}=0.
\]
\qstep{3}{Find the largest allowed acceleration.}
The meeting times are
\[
t_\pm=\frac{v_{\rm P}\pm\sqrt{v_{\rm P}^2-2av_{\rm P}t_{\rm w}}}{a}.
\]
For positive $a,v_{\rm P},t_{\rm w}$, real meetings after the wait require
$v_{\rm P}^2-2av_{\rm P}t_{\rm w}\geq0$. Therefore
\[
a_{\max}=\frac{v_{\rm P}}{2t_{\rm w}},\qquad
0<a\leq\frac{v_{\rm P}}{2t_{\rm w}}.
\]
At the boundary the roots merge and $t_*=2t_{\rm w}$.
\qstep{4}{Calculate the maximum.}
\[
a_{\max}=
\frac{\qty{12}{\metre\per\second}}{2(\qty{10}{\second})}
=\qty{0.600}{\metre\per\second\squared}.
\]
The meeting is at 20 s. Below the threshold, the first meeting is $t_-$;
a later meeting at $t_+$ occurs only if both continue.

\textit{Alternative:} maximize
$a(t)=\frac{2v_{\rm P}(t-t_{\rm w})}{t^2}$ for $t>t_{\rm w}$.
Its derivative $\frac{2v_{\rm P}(2t_{\rm w}-t)}{t^3}$ changes from positive to
negative at $t=2t_{\rm w}$, giving the same maximum.
\end{solution}

\quizpairbreak
\begin{problem}[25 pts]
\statementfigure{%
Two balls were thrown from the ground at angles of \ang{30} and
\ang{45} to the horizon. The ranges $R$ of the balls were the same.
Find the ratio of the initial speeds the balls had at the beginning
of their flights. Neglect air resistance.
}{\resizebox{\linewidth}{!}{\begin{tikzpicture}[font=\small]
\draw[physics line,-{Latex}] (-.15,0)--(5.6,0);
\draw[physics line,dashed,domain=0:5,samples=40]
 plot (\x,{\x*(1-\x/5)});
\draw[physics accent,dashed,domain=0:5,samples=40]
 plot (\x,{tan(30)*\x*(1-\x/5)});
\draw (0.65,0) arc[start angle=0,end angle=30,radius=.65];
\draw (1.05,0) arc[start angle=0,end angle=45,radius=1.05];
\node[fill=white,inner sep=1pt] at (.8,.18) {$30^\circ$};
\node[fill=white,inner sep=1pt] at (.8,.62) {$45^\circ$};
\node[below] at (0,0) {$0$};
\node[below] at (5,0) {$R$};
\node at (2.8,-.7) {Schematic, not to scale};
\end{tikzpicture}
}}
\end{problem}
\begin{solution}
\qstep{1}{Identify what is shared.}
The diagram shows equal launch and landing levels and a common range.
Use uniform downward gravity. The initial speeds are not assumed equal.
\sourcefig{0.48\linewidth}{projectile-speeds}
\qstep{2}{Relate speed to range.}
The nonzero flight time follows from vertical motion:
\[
t_f=\frac{2v_\theta\sin\theta}{g},\qquad
R=v_\theta\cos\theta\,t_f
=\frac{v_\theta^2\sin(2\theta)}{g}.
\]
\qstep{3}{Eliminate the common range and gravity.}
For two angles $\theta_A,\theta_B$,
$v_A^2\sin(2\theta_A)=v_B^2\sin(2\theta_B)$.
Taking the positive root because speeds are magnitudes,
\[
\frac{v_A}{v_B}
=\sqrt{\frac{\sin(2\theta_B)}{\sin(2\theta_A)}}.
\]
\qstep{4}{State the order and evaluate.}
Choosing the \ang{30} speed as the numerator gives
\[
\frac{v_{30}}{v_{45}}
=\sqrt{\frac{\sin90^\circ}{\sin60^\circ}}
=\sqrt{\frac{2}{\sqrt3}}\approx1.07.
\]
The reverse order is equally valid when labeled:
$\frac{v_{45}}{v_{30}}=\sqrt{\frac{\sqrt3}{2}}\approx0.931$.
Both are dimensionless. The \ang{30} launch needs the greater speed.

\textit{Alternative:} eliminate the two flight times directly from
$R=v_\theta\cos\theta\,t_f$ and
$0=v_\theta\sin\theta\,t_f-\frac{gt_f^2}{2}$; the same ratio follows.
\end{solution}

\quiznextproblem
\begin{problem}[25 pts]
A particle moves in uniform circular motion over a horizontal $xy$ plane
with the period of \qty{25}{\second}. At one instant, it moves through the
point with coordinates $[\qty{-1}{\metre},\qty{3}{\metre}]$ with the
velocity directed along $\hat{\mathbf i}$ and the acceleration of
$(0\hat{\mathbf i}+6\hat{\mathbf j})\,\unit{\metre\per\second\squared}$.
Find the $x$ and $y$ coordinates of the center of the circular path.
\end{problem}
\begin{solution}
\qstep{1}{Use the direction of centripetal acceleration.}
\sourcefig{0.37\linewidth}{circle-v2}
Centripetal acceleration points toward the center.
Here it is vertically upward, so $x_C=x$ and $y_C=y+R$.
\qstep{2}{Find the radius from centripetal acceleration.}
Let $a_c=|\mathbf a|$ be its magnitude and $v$ the constant speed.
One revolution covers $2\pi R$ in time $T$, so
\[
a_c=\frac{v^2}{R},\qquad v=\frac{2\pi R}{T}
\quad\Longrightarrow\quad a_c=\frac{4\pi^2R}{T^2}.
\]
\qstep{3}{Write the symbolic answer.}
\[
R=\frac{a_cT^2}{4\pi^2},\qquad
x_C=x,\qquad y_C=y+\frac{a_cT^2}{4\pi^2}.
\]
\qstep{4}{Calculate the radius and coordinates.}
\[
R=\frac{6(25)^2}{4\pi^2}\,\unit{\metre}
=\frac{1875}{2\pi^2}\,\unit{\metre}.
\]
\[
x_C=\qty{-1}{\metre},\qquad
y_C=\left(3+\frac{1875}{2\pi^2}\right)\unit{\metre}.
\]
\textbf{Answer:} $(x_C,y_C)\approx(-1,98.0)\,\unit{\metre}$.
The radius is added to the point's height; it is not the center's height itself.

\textit{Vector alternative:} with $\omega=\frac{2\pi}{T}$,
$\mathbf a=\omega^2(\mathbf r_C-\mathbf r)$ gives
$\mathbf r_C=\mathbf r+\frac{T^2}{4\pi^2}\mathbf a$, the same result.
\end{solution}
